% Motor: XeLaTeX. Sesión 22 - material docente UTP.
% Estructura docente de la Sesión 21: 20 diapositivas.
\documentclass[aspectratio=169,10pt,c]{beamer}
% Reutiliza sin cambios la identidad visual UTP validada en la Sesion 08.
% Reutiliza sin cambios la identidad visual UTP validada en la Sesion 05.
% La segunda ruta permite compartir los assets comunes sin duplicar binarios.
% !TEX root = slides.tex
% Preambulo comun: estilo UTP, colores, fuentes, plantillas de beamer y
% comandos de apoyo didactico. Se carga desde slides.tex tras \documentclass.

% La opción "c" de beamer reparte el espacio libre 1:1.5 (arriba:abajo), por lo
% que el contenido queda por encima del centro. Igualamos ambos rellenos.
\makeatletter
\define@key{beamerframe}{c}[true]{%
  \beamer@frametopskip=0pt plus 1fill\relax%
  \beamer@framebottomskip=0pt plus 1fill\relax%
  \beamer@frametopskipautobreak=0pt plus .5\paperheight\relax%
  \beamer@framebottomskipautobreak=0pt plus .5\paperheight\relax%
}
\makeatother

\usepackage{graphicx}
\usepackage{tikz}
\usepackage{xcolor}
\usepackage{iftex}

\ifPDFTeX
  \usepackage[utf8]{inputenc}
  \usepackage[T1]{fontenc}
  \usepackage{lmodern}          % tipografia escalable para el respaldo con pdflatex
  \renewcommand{\familydefault}{\sfdefault}
\else
  \usepackage{fontspec}
  \IfFontExistsTF{Arial}
    {\setsansfont{Arial}[BoldFont={Arial Bold}]}
    {\setsansfont{Helvetica}[BoldFont={Helvetica Bold}]}
\fi

\usepackage[spanish,es-nodecimaldot]{babel}
\usepackage{microtype}
\usepackage{booktabs}
\usepackage{array}
\usetikzlibrary{arrows.meta,positioning,fit,calc,shapes.geometric}

\graphicspath{{assets/}}

% ---------- Identidad visual UTP ----------
\definecolor{utpRojo}{HTML}{C42B35}
\definecolor{utpNegro}{HTML}{111111}
\definecolor{utpGris}{HTML}{5F5F5F}
\definecolor{utpGrisMedio}{HTML}{9A9A9A}
\definecolor{utpGrisClaro}{HTML}{F1F1F1}

\usetheme{default}
\setbeamercolor{normal text}{fg=utpNegro,bg=white}
\setbeamercolor{structure}{fg=utpRojo}
\setbeamercolor{frametitle}{fg=utpNegro,bg=white}
\setbeamercolor{title}{fg=utpNegro}
\setbeamercolor{subtitle}{fg=utpNegro}
\setbeamercolor{author}{fg=utpNegro}
\setbeamercolor{date}{fg=utpNegro}
\setbeamercolor{block title}{fg=white,bg=utpNegro}
\setbeamercolor{block body}{fg=utpNegro,bg=utpGrisClaro}
\setbeamercolor{alertblock title}{fg=white,bg=utpRojo}
\setbeamercolor{alertblock body}{fg=utpNegro,bg=utpRojo!7}
\setbeamercolor{exampleblock title}{fg=white,bg=utpGris}
\setbeamercolor{exampleblock body}{fg=utpNegro,bg=utpGris!8}
\setbeamercolor{terminal}{fg=white,bg=utpNegro}

\setbeamerfont{title}{series=\bfseries,size=\fontsize{19.5}{22.5}\selectfont}
\setbeamerfont{subtitle}{size=\fontsize{11.5}{13.5}\selectfont}
\setbeamerfont{date}{series=\bfseries,size=\fontsize{11.5}{13.5}\selectfont}
\setbeamerfont{author}{size=\fontsize{10.5}{12.5}\selectfont}
\setbeamerfont{frametitle}{series=\bfseries,size=\large}
\setbeamerfont{block title}{series=\bfseries,size=\small}
\setbeamerfont{block body}{size=\normalsize}

\setbeamertemplate{navigation symbols}{}
\setbeamertemplate{headline}{}
\setbeamertemplate{blocks}[default]
\setbeamertemplate{itemize item}{\color{utpRojo}\small\raise.15ex\hbox{$\blacktriangleright$}}
\setbeamertemplate{itemize subitem}{\color{utpGris}\scriptsize$\bullet$}
\setbeamertemplate{enumerate item}{\color{utpRojo}\insertenumlabel.}
\setbeamertemplate{background canvas}{\color{white}\rule{\paperwidth}{\paperheight}}

% Título de diapositiva: negrita negra a la izquierda, logo UTP a la derecha
\setbeamertemplate{frametitle}{%
  \begin{tikzpicture}[remember picture,overlay]
    \node[anchor=north east,inner sep=0pt]
      at ([xshift=-7mm,yshift=-5.5mm]current page.north east)
      {\includegraphics[width=17mm]{utp-logo.png}};
  \end{tikzpicture}%
  \vspace{1.2mm}%
  \begin{minipage}[t]{.80\textwidth}\raggedright\insertframetitle\end{minipage}\par
  \vspace{1.1mm}{\color{utpRojo}\rule{\textwidth}{1.1pt}}
}

% Solo numeración discreta; la plantilla UTP no lleva pie de página
\setbeamertemplate{footline}{%
  \hfill{\color{utpGrisMedio}\scriptsize\insertframenumber}\hspace{6mm}\vspace{2.5mm}%
}

% Portada UTP
\setbeamertemplate{title page}{%
  \begin{tikzpicture}[remember picture,overlay]
    \node[anchor=center,inner sep=0pt] at (current page.center) {%
      \includegraphics[width=\paperwidth,height=\paperheight]{utp-cover-bg.jpeg}%
    };
    \node[anchor=north,align=center,text=utpNegro,inner sep=0pt,
      text width=145mm,font=\usebeamerfont{title}]
      at ([yshift=-25.2mm]current page.north) {\inserttitle};
    \node[anchor=north,align=center,text=utpNegro,inner sep=0pt,
      text width=142mm,font=\usebeamerfont{subtitle}]
      at ([yshift=-38.2mm]current page.north) {\insertsubtitle};
    \node[anchor=north,align=center,text=utpNegro,inner sep=0pt,
      font=\usebeamerfont{date}]
      at ([yshift=-50.4mm]current page.north) {\insertdate};
    \node[anchor=north,align=center,text=utpNegro,inner sep=0pt,
      font=\usebeamerfont{author}]
      at ([yshift=-58.4mm]current page.north) {\insertauthor};
  \end{tikzpicture}%
}

% ---------- Comandos de apoyo didáctico ----------
\newcommand{\concepto}[2]{%
  \begin{block}{#1}\large #2\end{block}}
\newcommand{\pregunta}[1]{%
  \begin{alertblock}{Decide y explica}#1\end{alertblock}}
\newcommand{\analogia}[1]{%
  \begin{exampleblock}{En palabras simples}\small #1\end{exampleblock}}
\newcommand{\enlarealidad}[1]{%
  \begin{block}{Cómo se ve en la realidad}\small #1\end{block}}
\newcommand{\ojo}[1]{%
  \begin{alertblock}{Ojo con esto}\small #1\end{alertblock}}
\newcommand{\apoyo}{{\normalfont\footnotesize\color{utpGris}~· material de apoyo}}
\newcommand{\marca}[1]{\tikz[baseline=(m.base)]{%
  \node[rounded corners=1.5pt,fill=utpGrisClaro,draw=utpGrisMedio,inner xsep=3pt,inner ysep=1.5pt,
        font=\footnotesize] (m) {#1};}}
\newenvironment{terminal}{%
  \begin{beamercolorbox}[wd=\linewidth,sep=5pt,rounded=true]{terminal}%
  \ttfamily\small\ignorespaces}{%
  \end{beamercolorbox}}

\tikzset{
  caja/.style={draw=utpNegro, very thick, rounded corners=1.5pt, fill=white,
    align=center, minimum height=0.72cm, inner sep=5pt, font=\small},
  cajaazul/.style={caja, draw=utpGris, fill=white},
  cajaverde/.style={caja, draw=utpNegro, fill=utpGrisClaro},
  cajanaranja/.style={caja, draw=utpRojo!75, fill=utpRojo!7},
  flecha/.style={-{Stealth[length=2.3mm]}, very thick, draw=utpNegro},
  flechavuelta/.style={{Stealth[length=2.3mm]}-{Stealth[length=2.3mm]}, very thick, draw=utpNegro}
}

\graphicspath{{../../Sesion05/slides/assets/}}



\graphicspath{{assets/}{../../Sesion05/slides/assets/}}

\title{Integración de Servicios Cloud GCP}
\subtitle{IAM: quién puede hacer qué y sobre qué recurso}
\date{Semana 11 - Sesión 22}
\author{Rel Guzman}

\newcommand{\separador}[2]{%
\begin{frame}[plain]
\vfill
\begin{center}
{\fontsize{29}{34}\selectfont\bfseries #1}\par
\vspace{2mm}
{\small #2}\par
\vspace{4mm}
{\color{utpRojo}\rule{8.5cm}{1.2pt}}
\end{center}
\vfill
\end{frame}}

\begin{document}

{\setbeamertemplate{footline}{}\maketitle}

% 2
\begin{frame}{Revisión: la arquitectura de la sesión anterior}
\renewcommand{\arraystretch}{1.28}
\centering\small
\begin{tabular}{@{}p{3.0cm}p{4.9cm}p{5.1cm}@{}}
\toprule
\textbf{Bloque} & \textbf{Lo que decidimos} & \textbf{La pregunta que faltó}\\
\midrule
Backend en una VM & Ejecuta las reglas de la aplicación & ¿Con qué identidad guarda los archivos?\\
Bucket privado & Guarda facturas y reportes & ¿Quién puede leerlos y quién puede crearlos?\\
Proyecto de Google Cloud & Agrupa la VM, el bucket y la red & ¿Quién puede cambiar sus recursos?\\
\bottomrule
\end{tabular}
\par\vspace{3mm}
\begin{alertblock}{La pregunta de hoy}
\small En la Sesión 21 dijimos que el backend solo debe acceder a lo que necesita. Hoy vemos \textbf{cómo} se configura eso en Google Cloud: con IAM.
\end{alertblock}
\end{frame}

% 3
\begin{frame}{Revisión: una aplicación de facturas}
\begin{block}{Pedidos Express (proyecto \texttt{pedidos-express-prod})}
\small Laura emite facturas desde el navegador. La aplicación corre en una VM de Compute Engine, crea el PDF y el XML y los guarda en el bucket \texttt{pe-facturas-prod}.
\end{block}
\begin{center}
\begin{tikzpicture}[font=\small,
 caja/.style={draw=utpRojo,rounded corners=4pt,minimum width=2.55cm,minimum height=1.0cm,align=center}]
  \node[caja,fill=utpRojo!25] (u) at (0,0) {\textbf{Laura}\\\footnotesize navegador};
  \node[caja,fill=white,very thick] (v) at (5.2,0) {\textbf{VM}\\\footnotesize app de facturas};
  \node[caja,fill=utpRojo!10] (b) at (10.4,0) {\textbf{Bucket}\\\footnotesize PDF y XML};
  \draw[-latex,thick] (u) -- node[above,font=\footnotesize] {emite la factura} (v);
  \draw[-latex,thick] (v) -- node[above,font=\footnotesize] {guarda archivos} (b);
\end{tikzpicture}
\end{center}
\begin{alertblock}{Pregunta rápida}
\small Para que Laura pudiera descargar las facturas le dieron el rol \textbf{Editor} en todo el proyecto. Funciona. ¿Qué podría salir mal?
\end{alertblock}
\end{frame}

% 4
\begin{frame}[c]
\begin{block}{Logro de la sesión}
Al finalizar, podrás explicar cómo IAM controla el acceso en Google Cloud, asignar a cada identidad el \textbf{rol} mínimo sobre el \textbf{recurso} que necesita y usar una \textbf{cuenta de servicio} para que una VM acceda a Cloud Storage sin usar la cuenta de una persona.
\end{block}
\end{frame}

% 5
\begin{frame}{Saberes previos}
\begin{itemize}\setlength\itemsep{11pt}
  \item En una empresa, ¿la persona de contabilidad debería tener la llave del cuarto de servidores?
  \item Si una aplicación usa la cuenta personal de un desarrollador y esa persona deja la empresa, ¿qué le pasa a la aplicación?
  \item En el Lab 01 creaste un bucket. ¿Quién más podía ver sus archivos?
\end{itemize}
\end{frame}

% 6
\begin{frame}[c]{Utilidad de esta sesión}
\begin{center}
\begin{minipage}{.84\textwidth}
{\large En un proyecto real casi nunca trabajas solo. Contabilidad necesita leer facturas, un desarrollador necesita desplegar una VM y la aplicación necesita guardar archivos. Con IAM das a cada uno \textbf{solo lo que su tarea requiere}, y un error de una persona no derriba todo el proyecto.}
\end{minipage}
\end{center}
\end{frame}

% 7
\begin{frame}{Plan de hoy: dos partes y una práctica}
\renewcommand{\arraystretch}{1.35}
\centering\small
\begin{tabular}{@{}p{1.55cm}>{\raggedright\arraybackslash}p{5.8cm}>{\raggedright\arraybackslash}p{5.45cm}@{}}
\toprule
\textbf{Parte} & \textbf{Qué vemos} & \textbf{Lo conectamos con}\\
\midrule
1. Teoría & Identidad, rol y recurso: las tres piezas de IAM & El caso de facturas de Pedidos Express\\
2. Teoría & Cuentas de servicio y el rol \texttt{Service Account User} & El laboratorio \textit{Explora IAM}\\
\textbf{Práctica} & Corregir los permisos de un proyecto & Una tabla de identidad, recurso y rol\\
\bottomrule
\end{tabular}
\par\vspace{3mm}
\begin{alertblock}{Regla de hoy}
\small Si una identidad no necesita un permiso para su tarea, no se lo des.
\end{alertblock}
\end{frame}

% 8
\separador{Parte 1: quién, qué y sobre qué}{Las tres piezas de una asignación de IAM}

% 9
\begin{frame}{IAM en una frase}
\begin{center}
\begin{tikzpicture}[font=\small,
 caja/.style={draw=utpRojo,rounded corners=4pt,minimum width=3.0cm,minimum height=1.1cm,align=center}]
  \node[caja,fill=utpRojo!25] (q) at (0,0) {\textbf{Quién}\\\footnotesize identidad};
  \node[caja,fill=white,very thick] (r) at (4.6,0) {\textbf{Qué puede hacer}\\\footnotesize rol};
  \node[caja,fill=utpRojo!10] (s) at (9.2,0) {\textbf{Sobre qué}\\\footnotesize recurso};
  \draw[-latex,thick] (q) -- node[above,font=\footnotesize] {recibe} (r);
  \draw[-latex,thick] (r) -- node[above,font=\footnotesize] {en} (s);
\end{tikzpicture}
\end{center}
\renewcommand{\arraystretch}{1.24}
\centering\small
\begin{tabular}{@{}p{2.8cm}p{9.4cm}@{}}
\toprule
\textbf{Pieza} & \textbf{En el caso de facturas}\\
\midrule
Identidad & \texttt{laura@pedidosexpress.com}\\
Rol & \texttt{Storage Object Viewer}: listar y descargar objetos\\
Recurso & El bucket \texttt{pe-facturas-prod}\\
\bottomrule
\end{tabular}
\par\vspace{3mm}
\begin{block}{Se lee así}
\small ``Laura puede ver los objetos del bucket de facturas''. Nada más: ni borrarlos, ni tocar la VM.
\end{block}
\end{frame}

% 10
\begin{frame}{Identidades: personas y cuentas de servicio}
\begin{columns}[T,onlytextwidth]
\column{.48\textwidth}
\begin{block}{Cuenta de una persona}
\small Un correo como \texttt{laura@pedidosexpress.com}. La persona inicia sesión en la consola o en Cloud Shell.
\end{block}
\begin{exampleblock}{En el lab}
\small Trabajaste con dos usuarios: \textbf{Username 1} administraba el proyecto y \textbf{Username 2} solo tenía acceso de lectura.
\end{exampleblock}
\column{.48\textwidth}
\begin{block}{Cuenta de servicio}
\small Una identidad para una aplicación, no para una persona: \\\texttt{\footnotesize sa-generador-facturas@\\pedidos-express-prod.iam.gserviceaccount.com}
\end{block}
\begin{alertblock}{No inicia sesión}
\small Nadie escribe su contraseña. La VM la usa para llamar a otros servicios, como Cloud Storage.
\end{alertblock}
\end{columns}
\end{frame}

% 11
\begin{frame}{Roles: un paquete de permisos}
\renewcommand{\arraystretch}{1.3}
\centering\small
\begin{tabular}{@{}p{4.3cm}p{4.6cm}p{3.6cm}@{}}
\toprule
\textbf{Rol} & \textbf{Qué permite} & \textbf{Alcance}\\
\midrule
\texttt{Owner} / \texttt{Editor} / \texttt{Viewer} & Administrar, cambiar o ver casi todo & Muy amplio: roles básicos\\
\texttt{Storage Object Viewer} & Listar y descargar objetos & Solo Cloud Storage\\
\texttt{Storage Object Creator} & Crear objetos, sin leer ni borrar & Solo Cloud Storage\\
\texttt{Compute Instance Admin} & Crear y administrar VMs & Solo Compute Engine\\
\texttt{Service Account User} & Asociar una cuenta de servicio a una VM & Una cuenta de servicio\\
\bottomrule
\end{tabular}
\par\vspace{3mm}
\begin{alertblock}{La idea clave}
\small No asignas permisos sueltos: asignas un \textbf{rol}. Prefiere el rol más específico que cubra la tarea antes que \texttt{Editor}.
\end{alertblock}
\end{frame}

% 12
\begin{frame}{Dónde se asigna: proyecto o recurso}
\begin{center}
\begin{tikzpicture}[font=\small,
 caja/.style={draw=utpRojo,rounded corners=4pt,minimum width=2.6cm,minimum height=1.0cm,align=center}]
  \node[caja,fill=utpRojo!25,minimum width=3.4cm] (p) at (0,0) {\textbf{Proyecto}\\\footnotesize pedidos-express-prod};
  \node[caja,fill=white,very thick] (v) at (5.2,0.75) {\textbf{VM}\\\footnotesize facturación};
  \node[caja,fill=utpRojo!10] (b) at (5.2,-0.75) {\textbf{Bucket}\\\footnotesize pe-facturas-prod};
  \draw[-latex,thick] (p) -- (v);
  \draw[-latex,thick] (p) -- (b);
\end{tikzpicture}
\end{center}
\begin{columns}[T,onlytextwidth]
\column{.48\textwidth}
\begin{block}{Rol en el proyecto}
\small Vale para \textbf{todos} sus recursos. Se asigna en \textbf{IAM y administración > IAM}. En el lab diste \texttt{Storage Object Viewer} a Username 2 en el proyecto.
\end{block}
\column{.48\textwidth}
\begin{block}{Rol en un recurso}
\small Vale solo para \textbf{ese} recurso. Se asigna en la pestaña \textbf{Permisos} del bucket o de la cuenta de servicio.
\end{block}
\end{columns}
\par\vspace{2mm}
\begin{alertblock}{Elige el lugar más pequeño}
\small Si Laura solo usa un bucket, el rol va en ese bucket, no en el proyecto.
\end{alertblock}
\end{frame}

% 13
\begin{frame}{Mínimo privilegio: el problema de Laura}
\begin{columns}[T,onlytextwidth]
\column{.48\textwidth}
\begin{alertblock}{Antes: \texttt{Editor} en el proyecto}
\small Laura descarga facturas, pero también podría borrar la VM, cambiar la red o modificar otros recursos por error.
\end{alertblock}
\column{.48\textwidth}
\begin{exampleblock}{Después: \texttt{Storage Object Viewer} en el bucket}
\small Laura descarga las facturas para conciliarlas. No puede borrarlas ni tocar la infraestructura.
\end{exampleblock}
\end{columns}
\par\vspace{4mm}
\begin{block}{Principio de mínimo privilegio}
\small Cada identidad recibe el permiso mínimo necesario para realizar su tarea, sobre el recurso más pequeño posible. Si luego necesita más, se le concede solo esa acción.
\end{block}
\par\vspace{2mm}
\begin{block}{Dos pasos en la consola}
\small 1. \textbf{IAM y administración > IAM}: quitar a Laura el rol \texttt{Editor}.\\
2. \textbf{Cloud Storage > Buckets > pe-facturas-prod > Permisos}: agregar a Laura con \texttt{Storage Object Viewer}.
\end{block}
\end{frame}

% 14
\separador{Parte 2: cuentas de servicio}{Una VM que accede a Cloud Storage con su propia identidad}

% 15
\begin{frame}{Lo que hiciste en \textit{Explora IAM}}
\renewcommand{\arraystretch}{1.24}
\centering\small
\begin{tabular}{@{}p{4.6cm}p{4.3cm}p{3.6cm}@{}}
\toprule
\textbf{Paso} & \textbf{Qué hiciste} & \textbf{Qué comprobaste}\\
\midrule
Explorar la consola de IAM & Viste los roles de Username 1 y 2 & Quién tiene qué rol\\
Preparar un bucket & Username 1 creó un bucket y subió \texttt{sample.txt} & Username 2 podía verlo\\
Quitar acceso al proyecto & Retiraste el rol \texttt{Viewer} a Username 2 & Ya no ve el proyecto\\
Agregar permisos de almacenamiento & Diste \texttt{Storage Object Viewer} a Username 2 & Lista el bucket desde Cloud Shell\\
Configurar una cuenta de servicio & Creaste \texttt{read-bucket-objects} y la VM \texttt{demoiam} & La VM lee, pero no escribe\\
\bottomrule
\end{tabular}
\par\vspace{3mm}
\begin{alertblock}{Lo que muestra el lab}
\small Al quitar un rol amplio y dar uno específico, Username 2 conserva justo lo que necesita: leer objetos.
\end{alertblock}
\end{frame}

% 16
\begin{frame}{Una VM con cuenta de servicio}
\begin{center}
\begin{tikzpicture}[font=\small,
 caja/.style={draw=utpRojo,rounded corners=4pt,minimum width=2.7cm,minimum height=1.0cm,align=center}]
  \node[caja,fill=utpRojo!25] (d) at (0,0) {\textbf{Diego}\\\footnotesize despliega};
  \node[caja,fill=white,very thick] (v) at (5.2,0) {\textbf{VM}\\\footnotesize sa-generador-facturas};
  \node[caja,fill=utpRojo!10] (b) at (10.4,0) {\textbf{Bucket}\\\footnotesize pe-facturas-prod};
  \draw[-latex,thick] (d) -- node[above,font=\scriptsize,align=center] {Service Account\\User} (v);
  \draw[-latex,thick] (v) -- node[above,font=\scriptsize,align=center] {Storage Object\\Creator} (b);
\end{tikzpicture}
\end{center}
\begin{columns}[T,onlytextwidth]
\column{.48\textwidth}
\begin{block}{La cuenta de servicio}
\small \texttt{sa-generador-facturas} tiene \texttt{Storage Object Creator} sobre el bucket: crea el CSV diario, pero no lee ni borra facturas anteriores.
\end{block}
\column{.48\textwidth}
\begin{block}{Diego}
\small Tiene \texttt{Service Account User} sobre esa cuenta: puede elegirla al crear la VM, sin usar sus permisos desde su propia cuenta.
\end{block}
\end{columns}
\par\vspace{2mm}
\begin{exampleblock}{Igual que \texttt{demoiam} en el lab}
\small Desde la VM, \texttt{gcloud compute instances list} falló y la copia hacia el bucket fue rechazada; descargar \texttt{sample.txt} sí funcionó. La VM solo hace lo que permite el rol de su cuenta de servicio.
\end{exampleblock}
\end{frame}

% 17
\begin{frame}{Caso: la configuración de Pedidos Express}
\renewcommand{\arraystretch}{1.3}
\centering\small
\begin{tabular}{@{}>{\raggedright\arraybackslash}p{3.3cm}>{\raggedright\arraybackslash}p{3.4cm}>{\raggedright\arraybackslash}p{2.9cm}>{\raggedright\arraybackslash}p{2.9cm}@{}}
\toprule
\textbf{Identidad} & \textbf{Recurso} & \textbf{Rol} & \textbf{Motivo}\\
\midrule
Laura & Bucket \texttt{pe-facturas-prod} & \texttt{Storage Object Viewer} & Descarga facturas para conciliarlas\\
\texttt{\footnotesize sa-generador-facturas} & Bucket \texttt{pe-facturas-prod} & \texttt{Storage Object Creator} & La VM sube el CSV diario\\
Diego & Cuenta de servicio \texttt{\footnotesize sa-generador-facturas} & \texttt{Service Account User} & Asocia la cuenta a la VM\\
\bottomrule
\end{tabular}
\par\vspace{3mm}
\begin{columns}[T,onlytextwidth]
\column{.48\textwidth}
\begin{block}{La factura no se guarda con la cuenta de Laura}
\small Laura emite la factura; la VM guarda el PDF y el XML con su cuenta de servicio.
\end{block}
\column{.48\textwidth}
\begin{exampleblock}{Si Diego deja la empresa}
\small La aplicación sigue funcionando: usa su propia cuenta de servicio, no la de Diego.
\end{exampleblock}
\end{columns}
\end{frame}

% 18
\begin{frame}{Práctica: corrige los permisos de un proyecto}
\begin{block}{En grupos de dos}
\small Una librería tiene el proyecto \texttt{libreria-prod}, una VM que genera reportes de ventas cada noche y el bucket \texttt{libreria-reportes}. Hoy \textbf{todos} tienen \texttt{Editor} en el proyecto.
\end{block}
\renewcommand{\arraystretch}{1.15}
\centering\small
\begin{tabular}{@{}p{3.6cm}p{8.6cm}@{}}
\toprule
\textbf{Identidad} & \textbf{Tarea que realmente hace}\\
\midrule
Carla, contadora & Descarga los reportes para revisar las ventas\\
VM de reportes & Sube un reporte nuevo cada noche\\
Andrés, desarrollador & Crea la VM y le asocia su cuenta de servicio\\
Sofía, practicante & Solo necesita revisar la configuración, sin cambiar nada\\
\bottomrule
\end{tabular}
\par\vspace{0.5mm}
\begin{alertblock}{Entregable}
\small Una tabla con identidad, recurso, rol y motivo. Indiquen qué rol \texttt{Editor} retirarían.
\end{alertblock}
\end{frame}

% 19
\begin{frame}{Cierre: preguntas para comprobar}
\begin{block}{Responde con tus palabras}
\small
1. ¿Cuáles son las tres piezas de una asignación de IAM?\\\medskip
2. ¿Por qué dar \texttt{Editor} en el proyecto a Laura era un riesgo?\\\medskip
3. ¿Qué diferencia hay entre asignar un rol en el proyecto y en el bucket?\\\medskip
4. ¿Por qué la VM usa una cuenta de servicio y no la cuenta de un desarrollador?\\\medskip
5. ¿Para qué sirve el rol \texttt{Service Account User}?\\\medskip
6. En el lab, ¿por qué \texttt{demoiam} pudo leer \texttt{sample.txt} pero no escribir en el bucket?
\end{block}
\begin{alertblock}{Idea final}
\small IAM responde siempre tres preguntas: quién, qué puede hacer y sobre qué recurso. La mejor respuesta es la más pequeña que permite trabajar.
\end{alertblock}
\end{frame}

% 20
\begin{frame}{Referencias}
\scriptsize
\begin{itemize}\setlength\itemsep{4pt}
  \item Laboratorio: \textit{Explora IAM}, curso \textit{Essential Google Cloud Infrastructure: Core Services}: \texttt{cloudskillsboost.google/focuses/19082}.
  \item Material base de esta sesión: \texttt{/Users/rgap/Desktop/labs/ejemplo-explora-iam.md}.
  \item Diagrama del sistema desplegado: \texttt{labs/diagramas/iam\_sistema\_de\_facturas\_desplegado.md}.
  \item Simulación visual del caso: \texttt{labs/proyecto-facturas-web}.
  \item Material previo del curso: Sesión 21, arquitectura inicial de una webapp.
  \item Google Cloud: IAM, cuentas de servicio, Cloud Storage y Compute Engine.
  \item Las respuestas de las preguntas y de la práctica están en \texttt{../respuestas.md}.
\end{itemize}
\end{frame}

\end{document}
